% =============================================================================
% Бүлэг 2: Системийн зохиомж ба хэрэгжүүлэлт
% =============================================================================

\chapter{Системийн зохиомж ба хэрэгжүүлэлт}
\label{Chapter2}
\pagecolor{white}

\section{Шаардлага ба ерөнхий архитектур}
\label{sec:arch}

Системд дараах шаардлагуудыг тавьсан:
\begin{enumerate}
  \item \textbf{Монгол хэл.} Кирилл, латин галиг (хоёр хэв маяг), англи, холимог асуултад хариулах.
  \item \textbf{Байгууллагын дотор ажиллах.} Нээлттэй эхийн загвар, өөрийн GPU сервер (vLLM \cite{kwon_2023_vllm}).
        Өгөгдөл гадагш гарахгүй.
  \item \textbf{Аюулгүй байдал.} Зөвхөн унших эрх, хувийн мэдээлэл задруулахгүй, салбарын эрхийн хязгаарыг сахих.
  \item \textbf{Найдвартай байдал.} Хариулж чадахгүй, хоёрдмол асуултыг илрүүлж татгалзах эсвэл тодруулах.
  \item \textbf{Дахин давтагдах туршилт.} Бүх тохиргоо нэг YAML файлд (\tsql{config/experiments.yaml}) тодорхойлогдож, run бүрийн мета (загварын
        хувилбар, prompt, eval set, DB-ийн fingerprint, кодын хувилбар) бичигдэх.
\end{enumerate}

Зураг~\ref{fig:architecture}-т системийн бүрэлдэхүүнийг үзүүлэв. Код нь Python (≥ 3.10) дээрх \tsql{mn\_text2sql}
багц (≈14 700 мөр), туршилтын тохиргоо нь YAML, сервер дээр ажиллуулах скриптүүд нь bash хэлбэртэй.

\begin{figure}[htbp]
\centering
\begin{tikzpicture}[node distance=0.35cm and 0.35cm, every node/.style={font=\scriptsize},
  box/.style={draw, rounded corners, align=center, minimum height=0.8cm, text width=2.8cm, fill=gray!8},
  db/.style={draw, cylinder, shape border rotate=90, aspect=0.15, align=center, minimum height=1.5cm,
             text width=2.6cm, fill=blue!5}]
  \node[box] (in) {Асуулт\\(кирилл / латин /\\англи)};
  \node[box, right=of in] (ic) {Intent шалгалт\\(PII, бичих, хоёрдмол;\\зөвхөн үйлчилгээнд)};
  \node[box, right=of ic] (ret) {Хайлт: тайлбар,\\утга, нэр томьёо,\\жишээ};
  \node[box, right=of ret] (gen) {SQL үүсгэх\\Arctic / sqrl (vLLM)};
  \node[box, below=1.1cm of gen] (guard) {SQL хамгаалалт\\(sqlglot, allow-list)};
  \node[box, left=of guard] (exec) {Read-only гүйцэтгэх\\(timeout, EXPLAIN)};
  \node[box, left=of exec] (agent) {Засвар, санал\\хураалт, agent\\(tool + critic)};
  \node[box, left=of agent] (out) {Хариулт + тайлбар\\эсвэл татгалзах};
  \node[db, below=0.7cm of exec] (pg) {PostgreSQL:\\t2s view, RLS};
  \node[db, above=0.6cm of ret] (kn) {Мэдлэг: semantic\\model, glossary, индекс};
  \draw[->] (in) -- (ic); \draw[->] (ic) -- (ret); \draw[->] (ret) -- (gen); \draw[->] (gen) -- (guard);
  \draw[->] (guard) -- (exec); \draw[->] (exec) -- (agent); \draw[->] (agent) -- (out);
  \draw[<->, dashed] (exec) -- (pg); \draw[<->, dashed] (ret) -- (kn);
  \draw[->] (agent.north) -- ++(0,0.5) -| node[pos=0.25, below, font=\tiny]{алдаа бол засвар} ([xshift=-0.7cm]gen.south);
\end{tikzpicture}
\caption{Системийн бүрэлдэхүүн ба өгөгдлийн урсгал (алдаа гарвал засвар нь SQL үүсгэх алхам руу буцна)}
\label{fig:architecture}
\end{figure}

\section{Өгөгдлийн сан ба text2sql давхарга}
\label{sec:db}

\textbf{Схем.} Өмнөх ажлын репозиторт (text2sql-thesis) байсан, тэр ажлын туршилтад ашиглагдаагүй ББСБ-ийн PostgreSQL
схемийг ашиглав. 13 бизнес схемд нийт 189 хүснэгт ба view бий: байгууллагын үндсэн мэдээлэл (салбар, ажилтан),
харилцагч, KYC, бүтээгдэхүүн, зээл, барьцаа, данс ба гүйлгээ, хадгаламж, барагдуулалт (collection), эрсдэл, нягтлан
бодох бүртгэл, тайлан, систем. Ачаалах үед илэрсэн алдааг зөвхөн хуулбар дээр засаж, \tsql{db/CHANGES.md}-д бүртгэсэн.
Монгол Улсын хуулиар ББСБ иргэн, хуулийн этгээдээс хадгаламж авахыг хориглодог (харин итгэлцлийн үйлчилгээ эрхэлж
болно) \cite{mn_2002_nbfi_law}. Эх схемийн хадгаламжийн (\tsql{deposit}) хэсэг энэ утгаараа бодит ББСБ-ээс зөрөх
боловч итгэлцэлтэй төстэй бүтээгдэхүүний жишээ болгон туршилтад хэвээр үлдээв.

\textbf{Синтетик өгөгдөл.} Бодит харилцагчийн мэдээлэл ашиглах боломжгүй тул Python-оор детерминист (seed 42)
синтетик өгөгдөл үүсгэв. «Өнөөдөр» нь 2026-06-30, түүх 2024-01-01-ээс эхэлнэ. Үндсэн сан: 3000 харилцагч, 2666 зээл,
212 566 гүйлгээ, 17 салбар, 158 ажилтан. Зээлийн симуляц нь хуваарь, төлөлт, хугацаа хэтрэлт, бүтцийн өөрчлөлт,
шүүхийн шат, хасалт, ерөнхий дэвтрийн бичилтийг хамарна. DB-ийн функцүүд \tsql{CURRENT\_DATE} ашигладаг тул «өнөөдөр»-ийн
утга ачаалсан өдрөөс хамаарахгүй байлгахын тулд өгөгдлийг Python-д тооцож ачаалсан.

\textbf{text2sql давхарга (\tsql{t2s} схем).} Загварт 189 хүснэгт ба view-ийг шууд үзүүлэхийн оронд 28 view бүхий
тусгай схем бүтээв. View бүр нэг бизнес объектод тохирно (\tsql{loans}, \tsql{loan\_payments}, \tsql{collection\_cases},
\tsql{deposits}, \tsql{gl\_entries} г.м.). Код бүрийн хажууд нэрийн багана нэмж (\tsql{status\_code} → \tsql{status\_name}),
олон валютын дүнг төгрөгт хөрвүүлсэн багана (\tsql{*\_mnt}) нэмж, хувийн мэдээллийн багануудыг (регистр, утас,
хаяг) хассан. Үлдсэн харилцагчийн нэрийг semantic model-д хувийн мэдээлэл гэж тэмдэглэсэн. Туслах функц нь
\tsql{t2s.business\_date()} ба \tsql{t2s.fx\_to\_mnt()}.

\textbf{Эрх.} Bot хэрэглэгч (\tsql{t2s\_bot}) зөвхөн \tsql{t2s} схемд SELECT эрхтэй. Холболт бүр read-only
транзакц, 20 секундийн statement timeout, 2 секундийн lock timeout, 1 GB түр файлын хязгаартай. Мөрийн түвшний
хамгаалалт (RLS) нь салбарын хэрэглэгчид зөвхөн өөрийн салбарын мөрийг харуулна
(\tsql{t2s\_branch\_demo} хэрэглэгчээр туршсан).

\textbf{Test-suite хуулбарууд.} EX-ийн санамсаргүй тохиролтыг шалгахын тулд өөр seed-тэй, 0.25 масштабтай 8 хуулбар
DB (\tsql{nbfi\_ts1..8}) үүсгэв. Хуулбар бүр өөр төрлийн «хавх» агуулна: NULL утгыг олшруулсан, дүнг бөөрөнхийлж тэнцүү
утга үүсгэсэн, хугацаа хэтрэлт ба огнооны хил (30/31, 90/91 хоног, сарын эхэн ба сүүл), зарим хүснэгтийг хоосолсон
(хасалт, AML, батлан даагч, бүтцийн өөрчлөлт, хугацаанаас өмнөх татан авалт), эдгээрийн хослол ба энгийн өөр seed.

\section{Мэдлэгийн давхарга ба хайлт}
\label{sec:knowledge}

\textbf{Semantic model} (\tsql{knowledge/semantic\_model.yaml}). 28 view, 425 баганын монгол ба англи тайлбар,
кирилл синоним, enum кодын монгол нэр (\tsql{MOBILE = Мобайл апп}), гадаад түлхүүр, хувийн мэдээллийн тэмдэглэгээ.
Багана, хүснэгтийн нэр англиар үлддэг. Харин тайлбар нь DDL-ийн comment хэлбэрээр prompt-д орно
(Эх код~\ref{lst:ddl}).

\begin{lstlisting}[language=SQLmn, basicstyle=\ttfamily\scriptsize, caption={B1-ийн тайлбартай DDL-ийн хэсэг (prompt-д орох хэлбэр, товчилсон)}, label={lst:ddl}]
CREATE TABLE loan_applications (
  application_id TEXT,  -- (*@Хүсэлтийн ID@*)
  application_no TEXT,  -- (*@Хүсэлтийн дугаар@*) (LA-2025-000123); e.g. LA-2024-000001, LA-2024-000002
  status_code TEXT,  -- (*@Хүсэлтийн төлөв;@*) DRAFT=(*@Ноорог@*), SUBMITTED=(*@Илгээгдсэн@*), ..., REJECTED=(*@Татгалзсан@*), ...
  channel TEXT,  -- (*@Хүсэлт ирсэн суваг;@*) BRANCH=(*@Салбар@*), MOBILE=(*@Мобайл апп@*), WEB=(*@Вэб@*), AGENT=(*@Агент@*), ...
  requested_amount NUMERIC,  -- (*@Хүссэн дүн@*)
  ...
\end{lstlisting}

\textbf{Glossary} (\tsql{knowledge/glossary.yaml}). Байгууллагын 45 нэр томьёо: тодорхойлолт, каноник SQL, хамаарах
хүснэгт. Үүнээс 7 нь хоёрдмол («идэвхтэй зээл» = идэвхтэй багц уу, эсвэл зөвхөн хэвийн төлөвтэй юу).
Үйлчилгээний горимд intent шалгалт хоёрдмол үгийг тодотголгүй асуувал тодруулах асуулт буцаана. Туршилтад энэ
шалгалт унтраалттай тул хоёрдмол асуултад загвар өөрөө хариулна (§\ref{sec:eval-method}).

\textbf{KPI-ийн semantic layer} (\tsql{knowledge/metrics.yaml}). 20 гол үзүүлэлт, 14 хэмжээс (185 зөвшөөрөгдсөн хос).
F8 туршилтад жижиг загвар нь асуултыг \{metric, хэмжээс, шүүлт\} JSON болгож, SQL-ийг кодоор compile хийнэ.

\textbf{Индекс ба хайлт.} \tsql{index-build} нь дөрвөн индекс үүсгэнэ:
(1) схемийн баримт;
(2) текст баганын давтагдашгүй утга (хувийн мэдээллийн баганаас бусад, ≤400 утгатай);
(3) нэр томьёо;
(4) 160 жишээ асуулт.
Хайлт нь BM25 ба BGE-M3 embedding-ийн оноог $\lambda = 0.5$ жингээр хослуулна. Утгын хайлт гурван горимтой:
\begin{itemize}
  \item \textbf{basic} (B2): кирилл түлхүүр ба галигийн «сул түлхүүр». Сул түлхүүр нь о/ө/у/ү-г нэг үсэг болгож,
        kh/h, ts/c, w/v-г нэгтгэж, давхар үсгийг дан болгоно. Латин «v» нь ү эсвэл в байж болох тул хоёр хувилбараар хайна;
  \item \textbf{rule} (F2): + дагавар таслах (-ийн, -аас, -д …) ба галигийн хоёр хэв маягийн хувиргалт;
  \item \textbf{learned} (F3): + хэрэглэгчийн лог ба засвараас сурсан галиг → кирилл хос.
\end{itemize}
Few-shot жишээг DAIL-SQL маягаар асуултын утгыг \tsql{<утга>}, \tsql{<тоо>}-оор далдалж төстэй байдлаар сонгоно
\cite{gao_2024_dail_sql}. Схемийн тайлбар токены төсвөөс (13 000) хэтэрвэл бүх хүснэгт DDL-д үлдэж, тайлбарыг зөвхөн
асуулттай хамгийн холбоотой хүснэгтүүдэд өгнө.

\section{Загварын протокол ба SQL хамгаалалт}
\label{sec:protocol}

Туршилтын хоёр гол загварыг \textbf{тус бүрийн сургалтын протоколоор} ажиллуулна. Учир нь RL-ээр сургасан загвар
өөр prompt хэлбэрт муудаж болох бөгөөд sqrl-9b-ийн карт ч сургалтын протоколоо яг мөрдөхийг зөвлөдөг
\cite{feyn_2026_sqrl}.

\textbf{Arctic-Text2SQL-R1-7B} (Qwen2.5-Coder-7B-Instruct-д суурилсан, OmniSQL-оос эхлүүлж GRPO хийсэн)
\cite{yao_2025_arctic_text2sql_r1, snowflake_2025_arctic_card}. A1-ийн prompt нь өгүүллийн Зураг C.1-тэй ижил (бусад
тохиргоонд жишээ ба нэмэлт зааврын хэсэг нэмэгдэнэ): system нь «data science expert», user нь
Database Engine / Database Schema / Question / Instructions / Output Format хэсэгтэй. Assistant-ийн хариуг
«Let me solve this step by step.\textbackslash n<think>» гэж эхлүүлж (prefill), SQL-ийг \tsql{<answer>} доторх
sql код блокоос авна. RAG-ийн агуулга нь DDL-ийн comment ба асуултын өмнөх evidence хэсэгт орно.

\textbf{sqrl-9b} (Qwen3.5-9B суурьтай agentic загвар) \cite{feyn_2026_sqrl}. \tsql{<schema>}, \tsql{<evidence>} хэсэгтэй
system prompt. Загвар \tsql{<sql>...</sql>}-ээр DB-г судалж (≤5 алхам), ажиглалтыг \tsql{<observation>}-оор
авна. Эцэст нь \tsql{<answer>}-ээр хариулна.

\textbf{Диалект.} Хоёр загвар голчлон SQLite дээр сургагдсан. Иймд гурван аргыг туршина:
\begin{enumerate}
  \item SQLite протоколоор шууд асууж, гарсан SQL-ийг PostgreSQL дээр гүйцэтгэх (A1);
  \item prompt-д PostgreSQL-ийг заах (A6);
  \item SQLite-аар үүсгээд sqlglot-оор PostgreSQL руу хөрвүүлэх (A7) \cite{mao_2026_sqlglot}.
\end{enumerate}

\textbf{SQL хамгаалалт.} Загварын гаргасан SQL бүр (sqrl-ийн судлах асуулга ч) дараах шалгалтыг давна:
\begin{enumerate}
  \item sqlglot-оор задлах; ганц \tsql{SELECT}/\tsql{WITH} асуулга байх;
  \item өгөгдөл өөрчлөх, DDL, \tsql{SET}, \tsql{COPY}, \tsql{SELECT INTO}, \tsql{FOR UPDATE} хориотой;
  \item хориотой функцүүд (\tsql{pg\_sleep}, \tsql{dblink}, \tsql{pg\_read\_file}, \tsql{set\_config} г.м. 35 гаруй);
  \item зөвхөн \tsql{t2s} схемийн бүртгэлтэй view (alias-тай баганыг allow-list-ээр, бусдыг EXPLAIN-ээр шалгана);
  \item \tsql{CURRENT\_DATE}, \tsql{now()}, \tsql{date('now')}, нэг аргументтай \tsql{age()}-ийг бизнесийн
        өдрөөр дахин бичих (үнэлгээ хуанлийн өдрөөс хамаарахгүй);
  \item EXPLAIN-ийн өртгийн дээд хязгаар ($5 \times 10^7$), мөрийн хязгаар.
\end{enumerate}
Энэ нь эхний шүүлтүүр бөгөөд DB-ийн эрх, RLS, read-only транзакц нь хоёр дахь, бие даасан хамгаалалт болно.

\textbf{Засвар, санал хураалт, татгалзалт.} E1 нь алдааны мессеж (эсвэл хоосон үр дүн) буцааж 2 удаа хүртэл засуулна
(төлөвлөгөөнд 1--3). Загвар SQL бичихээс татгалзсан (SQL-гүй бүрэн хариу өгсөн) бол засуулахгүй, «татгалзсан» гэж
тооцно. Тасарсан гаралт ба SQL олдоогүй форматын алдаа нь алдаа хэвээр бөгөөд E1 засуулна.
E2 нь 8 нэр дэвшигч SQL үүсгэж (sampling), гүйцэтгэлийн үр дүнгийн hash-аар санал хураана. Олонхын хувь нь итгэлийн
хэмжүүр болно. E3 нь итгэл босгоос доош бол хариулахаас татгалзана.

\section{Agentic урсгал}
\label{sec:agent}

Өмнөх ажлын Planner → Writer → Critic агентын санааг шинэ системийн хамгаалалт, хайлт дээр дахин зохиов
(Зураг~\ref{fig:agent}). Загвар өөрийгөө засах чадвар гадаад дохиогүй үед найдваргүй
\cite{huang_2024_self_correct}. Иймд agent-ийн гол дохио нь \textbf{детерминист tool}-ууд байна:
\begin{itemize}
  \item \textbf{Гүйцэтгэлийн шалгалт.} Алдаа гарвал заавал засуулах хатуу олдвор болно. Харин хоосон үр дүн ба бүгд NULL нь
        зөөлөн олдвор: зөв хариу хоосон байж болох тул шүүмжлэгчид л дамжуулна.
  \item \textbf{Литерал шалгах tool.} \tsql{WHERE}, \tsql{IN}, \tsql{LIKE} дахь текст утга тухайн баганад үнэхээр хадгалагдсан
        эсэхийг DB-ээс шалгана. Байхгүй бол тухайн баганын хамгийн ойр хадгалсан утгуудыг санал болгоно. Хувийн
        мэдээллийн баганын утгыг харуулахгүй. Үр дүн хоосон биш үед олдвор зөөлөн байна.
  \item \textbf{LLM шүүмжлэгч.} Асуулт, SQL, үр дүнгийн эхний 8 мөр, tool-ийн олдворыг өгөхөд загвар
        «VERDICT: OK» эсвэл «VERDICT: REVISE --- …» гэж дүгнэнэ. Дүгнэлтгүй хариуг SQL өөрчлөхгүйгээр
        тусад нь бүртгэнэ.
\end{itemize}
Засварласан SQL гүйцэтгэгдэхгүй, эсвэл ажиллаж байсан үр дүнг хоосон болговол түүнийг үгүйсгэж, өмнөх SQL-ийг
хадгална. Загвар SQL бичихээс татгалзсан үед (хариулах боломжгүй асуулт) agent SQL бичүүлэх гэж шахахгүй.
Сонголтоор planner нь SQL бичихээс өмнө англи төлөвлөгөө (хүснэгт, join, шүүлт, агрегат) гаргана.

\begin{figure}[htbp]
\centering
\begin{tikzpicture}[node distance=0.5cm and 0.6cm, every node/.style={font=\scriptsize},
  box/.style={draw, rounded corners, align=center, minimum height=0.8cm, text width=2.4cm, fill=gray!8},
  dec/.style={draw, diamond, aspect=2.4, align=center, inner sep=0pt, fill=yellow!10, font=\scriptsize}]
  \node[box] (plan) {Planner (G3)\\англи төлөвлөгөө};
  \node[box, right=of plan] (write) {Writer\\(загварын протокол)};
  \node[box, right=of write] (obs) {Tool: гүйцэтгэл,\\литерал шалгалт};
  \node[box, below=1.0cm of obs] (crit) {Critic\\VERDICT};
  \node[dec, left=0.9cm of crit] (d) {хатуу олдвор\\эсвэл REVISE?};
  \node[box, left=0.9cm of d] (fin) {Хариулт};
  \draw[->] (plan) -- (write); \draw[->] (write) -- (obs); \draw[->] (obs) -- (crit); \draw[->] (crit) -- (d);
  \draw[->] (d) -- node[above]{үгүй} (fin);
  \draw[->] (d.north) -- node[right, pos=0.45]{тийм (≤2)} (write.south);
\end{tikzpicture}
\caption{Agentic давталт (G бүлэг)}
\label{fig:agent}
\end{figure}

\section{Fine-tuning-ийн дамжлага}
\label{sec:ft-pipeline}

Сургалтын өгөгдөл гурван эх сурвалжтай:
\begin{enumerate}
  \item BIRD train-ээс DB-ээр давхаргалж сонгосон 5000 асуулт. Apache лицензтэй загвараар (Qwen3.5-9B) монгол руу
        орчуулна. Орчуулахаас өмнө асуултад гарсан SQL-ийн утга, нэрийг placeholder-оор далдалж, дараа нь сэргээнэ
        (BIRDTurk-ийн идентификатор ба литералыг хөндөхгүй байх зарчмаас санаа авсан \cite{aktas_2026_birdturk}).
        30\%-ийг галиг хувилбараар үүсгэнэ.
  \item Өөрийн DB-ийн хэв загвар асуултын тогтмолыг сольж үүсгэсэн синтетик асуулт (≤3000), гүйцэтгэлээр шалгасан.
  \item Жишээний сангийн 160 асуулт (\tsql{knowledge/examples.jsonl}). Eval set-тэй ижил аргаар бэлтгэсэн
        (§\ref{sec:evalset}), eval-ийн асуулттай давхцахгүй.
\end{enumerate}
Dev/test асуулттай ойролцоо (тэмдэгтийн n-gram ≥ 0.8 эсвэл embedding cosine ≥ 0.9) эсвэл gold-той ижил үр дүнтэй
мөрийг хасна.

Дамжлага (\tsql{t2s\_train}): \tsql{fetch → translate → synth → dedup → sample → cot → build → sft → grpo}.
\begin{itemize}
  \item \textbf{C1 (rejection sampling).} Асуулт бүрт 8 хариу ($T = 0.8$); гүйцэтгэлээр зөвийг ≤4 хадгална.
  \item \textbf{C2 (CoT).} Нэг ч зөв хариугүй асуултад gold-оор чиглүүлсэн reasoning үүсгэж, «gold-ийг хуулсан» гаралтыг хасна.
  \item \textbf{SFT.} LoRA $r = 32$, $\alpha = 64$, бүх linear давхарга, LR $2\times10^{-4}$, cosine, 2 epoch,
        bf16, loss зөвхөн completion дээр.
  \item \textbf{GRPO.} 2000 «заримдаа шийддэг» (8-аас 1--7 нь зөв) асуулт, $G = 8$, $\beta = 0.001$, LR $10^{-6}$,
        reward = 1 / 0.1 / 0 (§\ref{sec:finetune}), давуу талыг тэгшитгэл~\ref{eq:grpo}-ээр тооцно. TRL GRPOTrainer,
        vLLM colocate. TRL олон ээлжит tool орчин дэмждэггүй тул sqrl-9b-ийн GRPO нь нэг ээлжит (DB судлах алхамгүй)
        хэлбэрээр явагдана.
  \item \textbf{F4.} Нэг асуултын кирилл, галиг, англи хувилбарыг нэг batch-д оруулж, бусад хувилбарын олонхтой ижил
        үр дүн гаргасан rollout-д +0.2 нэмнэ (бичгээс үл хамаарах reward; хариу буруу ч олонхтой таарвал урамшуулна).
\end{itemize}
Сургалт нь A100 80 GB дээр явагдана. Үнэлгээг L4 (24 GB) дээр хийнэ: SFT-ийн LoRA adapter-ийг vLLM-д бүртгэж,
GRPO-ийн үр дүнг суурь жинтэй нэгтгэж (merge) тусдаа загвар болгон ажиллуулна. Туршилтын төлөвлөгөөний дагуу L4 дээр
Arctic-ийг BF16-аар, sqrl-9b-ийг санах ойд багтаахын тулд FP8 квантчиллаар ажиллуулна.

\section{Үйлчилгээ ба аюулгүй байдал}
\label{sec:service}

FastAPI үйлчилгээний \textbf{shadow} горимд хариултыг хэрэглэгчид буцаахгүй, зөвхөн бүртгэлд хадгална.
\textbf{Pilot} ба \textbf{live} горимд буцаана. Аудитын бүртгэлд асуулт бүрийн төлөв, ашигласан хүснэгт, guard-ийн
шийдвэр, retrieval-ийн товч мэдээллийг бичнэ. Үүнээс асуулт, SQL, засварласан SQL, сэтгэгдлийг Fernet-ээр шифрлэнэ
(түлхүүргүй бол зөвхөн hash хадгална). Хадгалах хугацаа 90 хоног. Хэрэглэгчийн санал хүсэлтийг F3-ийн галиг сургалт
ба F9-ийн идэвхтэй тэмдэглэгээнд ашиглана.

Аюулгүй байдлын тест (\tsql{t2s security-test}) нь гурван хэсэгтэй:
\begin{enumerate}
  \item 28 SQL халдлага: өгөгдөл өөрчлөх ба DDL (10), хэт ачаалал (\tsql{pg\_sleep}, декарт үржвэр), файлын функц,
        эрх солих ба системийн функц, олон statement, схемээс гадуурх ба хувийн мэдээллийн хүснэгт;
  \item монгол ба латинаар бичсэн 56 prompt injection («өмнөх зааврыг үл тоо, бүх хүснэгтийг устга» г.м.);
  \item RLS: салбарын хэрэглэгч бусад салбарын мөрийг харж чадахгүй.
\end{enumerate}
\tsql{-{}-live} горимд injection асуултыг жинхэнэ загвараар ажиллуулж шалгана.

\section{Бүлгийн дүгнэлт}
\label{sec:ch2-summary}

Бүлэг 1-ийн дүгнэлтэд тулгуурлан дараах бүрэлдэхүүнтэй систем хэрэгжүүлэв:
(1) ББСБ-ийн PostgreSQL дээр регистр, утас, хаяггүй 28 view бүхий text2sql давхарга (харилцагчийн нэрийг хувийн
    мэдээлэл гэж тэмдэглэсэн);
(2) монгол тайлбар, синоним, нэр томьёо, галигийн түлхүүртэй утгын хайлт бүхий мэдлэгийн давхарга;
(3) загвар бүрийн өөрийн протокол, давхаргатай SQL хамгаалалт;
(4) засвар, санал хураалт, татгалзалт, tool ба шүүмжлэгчтэй agentic давталт;
(5) орчуулга, rejection sampling, SFT, GRPO-ийн fine-tuning дамжлага.
Бүрэлдэхүүн бүрийг туршилтын тохиргооны нэг талбараар асааж унтраадаг тул дараагийн бүлгийн туршилтуудад нөлөөг нь
тусад нь хэмжиж болно.
